\section{A route can carry phase before it carries a record}
A detector need not be watched for it to change an experiment. A physical interaction can correlate a route with another degree of freedom, and that correlation can change the interference seen when the route is recombined. The further question is whether the correlation remains readable, by which apparatus, and for how long. These are physical questions about a carrier and its dynamics, not about an observer's awareness.

Unified Volume I calls the transition to a durable accessible record a substrate write.\cite{unified} Its quantum section retains the Born rule and action phase, while deferring the detailed write algebra to the computational volumes. The present chapter expands that interface. It does not infer an amplitude evolution from a route label or turn an accessibility convention into a new collapse law.

Chapters I--III provide an explicit interacting source, its clock current, and a source-to-field interface. Chapter IV separates local clock evolution from signal transport. Chapter V distinguishes a traveler's history from what a receiver can obtain and from what a prescribed boundary allows.\cite{chapter1,chapter2,chapter3,chapter4,chapter5} The quantum continuation asks what happens when alternatives remain coherent through part of that chain.

Three questions guide the calculation:
\begin{enumerate}[leftmargin=*,itemsep=2pt]
\item What phase and amplitude does a declared dynamics assign to each coherent alternative?
\item Which physical degrees of freedom become correlated with those alternatives, and how does that change interference?
\item What record can a specified domain actually decode, retain and transmit?
\end{enumerate}
A free photon is not unwritten because it is massless. It can carry energy, momentum and phase without yet leaving the kind of persistent record being discussed. Nor does an inaccessible correlation stop affecting a reduced quantum state. Accessibility concerns which records a domain can use; the full joint evolution still contains correlations outside that domain.

For an orthonormal set of exclusive route-readout alternatives, normalized amplitudes give
\begin{equation}
 P(j)=\frac{|a_j|^2}{\sum_k|a_k|^2}.
\end{equation}
This is the retained Born premise.\cite{Born1926} Coherent alternatives reaching the same complete final outcome contribute to a common amplitude before its squared magnitude is taken. Arbitrary decompositions of one state do not define additional mutually exclusive outcomes. After a record forms, subsequent and conditional measurements still obey the same probability rule.

\section{What the recovered source record establishes}
The recovered CR075 package reports a scoped structural ledger-write pass. Its runner inspects source text, manifests, disclosures and named bridge entries. The retained QP reports supply constructed route tables and threshold selectors. Their presence is evidence of a specified computational interpretation, not a derivation of arbitrary complex amplitudes or a physical detector dynamics.\cite{cr075}

This distinction matters for the probability account. QP014 normalizes constructed capacity, coherence-survival and access weights. It does not establish that these weights equal $|a_j|^2$ for arbitrary quantum preparations. QP001, QP010 and QP016 likewise do not establish universal physical decoherence or outcome thresholds. This chapter therefore uses their route/write distinction while deriving its interference calculations from explicitly stated quantum laws.

The later exact-write sources supply reversible finite address updates and separately retained history data. A signed update followed by its inverse can restore an address while a retained event list remains different. That is a useful exact algebraic example of restoration without deletion. It is not yet a microscopic measurement interaction.\cite{writealgebra}

The lifecycle source also keeps an important boundary open: the maps from its computational \texttt{A\_OPERATOR} to contact and to \texttt{X1} are not supplied as a completed physical law. That typed operator must not be silently identified with the scalar accumulation coordinate $A(r)$ or with a Hamiltonian coupling.\cite{lifecycle} The task here is to state what such a bridge would need to reproduce.

All recovered repository bytes are pinned and authenticated in the companion package. The complete CR075 directory was inspected; its original 242-input upstream environment was not recovered or rerun. Historical reported control flags are not recast as newly executed perturbed quantum experiments. The derivations and checks below are a new independent receiving calculation, not an upgrade of those historical receipts.

\section{Three phases and one consistent action account}
\subsection{A signed Write integral is not a complete history}
The original lifted source phases define
\begin{equation}
\begin{gathered}
\gamma_{\rm W}=\phi_a+\phi_b-2\phi_p,
 \\
 \chi=\phi_p,
\end{gathered}
\end{equation}
with the legal-path one-form
\begin{equation}
 \omega_{\rm Write}=\hbar(\dd\gamma_{\rm W}+3\dd\chi).
\end{equation}
Consequently its integrated signed action is
\begin{equation}
 S_{\rm Write}=\hbar(\Delta\gamma_{\rm W}+3\Delta\chi).
\end{equation}
The lifted endpoints, wrapped phases, ordered event word and elapsed duration preserve different information. Chapter I's retained histories can have equal endpoint and equal signed action while retaining different positive durations and event orders.\cite{chapter1} An exact differential does not reconstruct an ordered word from its integral.

The interacting phase $\theta=-\arg\langle Z\rangle$ is a separate observable coordinate. An interferometer's relative phase is another operand, obtained from its full controlled amplitude. No equality among $\gamma_{\rm W}$, $\theta$ and the observed fringe phase follows simply because all three are called phases. Their identification requires an explicit coupling and readout.

\subsection{An explicit massive-route comparison}
For a declared test-body action, the massive point contribution is
\begin{equation}
 S_{{\rm mass},j}=-m_0c^2\int_{\gamma_j}\dd\tau.
 \label{eq:massive-action}
\end{equation}
Here $m_0c^2$ must be a part of the energy not already counted in the internal generator. In the weak static comparison with $F=1-A$ to first order, $A\ll1$, and $v^2/c^2\ll1$,
\begin{equation}
 \frac{\dd\tau}{\dd t}
 =1-\frac A2-\frac{v^2}{2c^2}
    +O(A^2,Av^2/c^2,v^4/c^4).
\end{equation}
For two routes with the same coordinate duration, the common rest phase cancels and
\begin{equation}
 \Delta\varphi_{\rm mass}
 =\frac{m_0}{2\hbar}\int
       \left[c^2(A_1-A_0)+(v_1^2-v_0^2)\right]\dd t
       +\text{higher orders}.
\end{equation}
This is the ordinary phase of kinetic energy minus Newtonian potential energy with $\Phi=-c^2A/2$. It is conditional on the stated action and geometry, not a derivation of an $A$-dependent quantum law from the source grammar.

An interferometer also needs the phases of its preparation, splitting, routing and recombination operations, with recoil and every record retained by those controls. A comparison of two held clocks is not automatically a coherent recombination experiment. For static segments of duration $T$, the proper-age difference is $T(\sqrt{F_1}-\sqrt{F_0})$ under the declared metric; the apparatus must still make both alternatives and their interference available.

Equation~\eqref{eq:massive-action} is not a photon phase prescription: a null route has zero proper time. A photon requires its wave equation and boundary-value or eikonal phase, with the source and receiver controls included. The earlier photon-road delay does not by itself define a quantum-state amplitude.

